% 補足資料 S12 の表：事前登録の記録の一覧（第7弾 2026-10-28，第8弾 10-29 に組4〜6 の2行，第9弾 10-29 に結果のコミットの列と符号無作為条件の組1〜3 の結果を足した）。
% 出典：docs/ の各作業記録（段階I・II，照合，1ビット，追加実験 paper_extra_log.md）の事前登録の節と，git log -S でその文が最初に入ったコミット。
%   結果のコミットは git の題と日時で確かめた（paper/check/numbers_check_v7.md §3，numbers_check_v9.md §1）。
%   「事前」「事前登録外」の行は登録のコミットが結果のコミットより前（同じ日の行も時刻で前）。
%   案と確定が別のコミットの行は確定の方を書いた（案のコミットは numbers_check_v7.md §3）。
\setlength{\tabcolsep}{4pt}%
\begin{tabular}{@{}p{0.48\linewidth}p{0.15\linewidth}p{0.15\linewidth}p{0.12\linewidth}@{}}
\toprule
何を登録したか（判定・予想） & 登録 & 結果 & 区別 \\
\midrule
段階I（経験50\%）の他者側の関門 G5a・G5b & 10-19 ed26338 & 10-20 b727ca5 & 事前 \\
段階I（経験50\%）の関門の確定（G1 は前提，G3 は主，判定の点） & 10-19 c672c33 & 10-20 f1df9c8 & 事前 \\
関門の結果による段階IIの扱い & 10-20 aa919ba & ― & 事前 \\
組4 を予備とし除外しない（4組では3組以上） & 10-20 d0256e4 & ― & 事前 \\
段階II（経験50\%）の構成と T2・T5 の判定 & 10-20 99db5b9 & 10-21 0076911 & 事前 \\
形制限モジュール（自由出力の結果の後に設計，探索） & 10-21 5a51374 & 10-21 1b48d61 & 事後 \\
他者遮蔽学習（三層の分析から導いた仮説） & 10-21 75780e9 & 10-22 df28e39 & 事前（探索） \\
段階I・II（経験80\%）．予想「$g$ と G5a の傾きが経験50\%より大きい」 & 10-22 1088e06 & 10-23 6c03964，10-24 7815c6e & 事前 \\
照合（経験50\%）の手続き・候補・主の指標・判定 & 10-23 b0da18f & 10-23 7600eff & 事前 \\
1ビット（経験50\%），照合を経験80\%へ，好み切替条件 & 10-23 093e509 & 10-24 d004d47，0f762a4 & 事前 \\
他者特徴基準の数え方の訂正 & 10-23 8256ec1 & ― & 事後（訂正） \\
他者特徴基準（MSE 期間）と好み追従基準 & 10-23 ae2d995 & 10-24 d004d47 & 事前 \\
G5a を副に，G5b を他者側の主に & 10-23 d659924 & ― & 事後 \\
符号に従う分の定義（組1〜4 では事後） & 10-23 3737101 & 10-23 83ad20c & 事後 \\
1ビットの出力と A-2 の方角の比較 & 10-24 bf9df2c & 10-24 24a93f0 & 事後 \\
1ビット（経験80\%）と「好みを書く」条件 & 10-24 66194ad & 10-24 072c94d & 事前 \\
照合を組4（経験80\%）と非共有条件へ & 10-24 6dd4741 & 10-24 f900fb6，10-25 1613914 & 事前 \\
組5・6 の追加と符号に従う分の経験の比較 & 10-24 00b8e48 & 10-25 6709d65 & 事前 \\
照合を組5・6 へ & 10-26 4d167b3 & 10-26 e248bed & 事前 \\
他者の規則の確率を 0.8 にした比較（$\Delta$） & 10-26 29d824f & 10-26 e42e4e6 & 事前 \\
学習後に重みを全部写す比較 & 10-26 29d824f & 10-26 c4fba2e & 事前 \\
符号無作為条件の組1〜3：V1（G5b・符号に従う分），V2（$\Delta$），V3（$d_P$）．V1・V3 は成立，V2 は不支持 & 10-26 f61a3b3 & 10-27 02a7127 & 事前 \\
快・不快の列だけを写す条件 & 10-26 d7850b0 & 10-26 c4fba2e & 事前 \\
零入力条件 & 10-26 d7850b0 & 10-27 a39b3d4 & 事前 \\
経験なし条件のモデルに写す確認（測り方） & 10-27 fbe12c5 & 10-27 94bf4ea & 事前登録外 \\
追跡のできる評価データでの比較（組1〜3） & 10-28 fae2e8b & 10-28 60591a2 & 事前 \\
除外の影響の読み方（S5） & 10-28 3c6cd98 & 10-28 78bb877 & 事前登録外 \\
好み切替条件の $d_P$ の読み方 & 10-28 6d955dd & 10-28 0505fa9 & 事前登録外 \\
符号無作為条件の組4〜6 の追加と6組の判定（V2 の予想は置かず記述） & 10-28 bb6de41 & 10-29 91b0f7f & 事前 \\
値の取り方の追記（食い違えば弱い方で書く） & 10-28 b12683d & 10-29 91b0f7f & 事前 \\
\bottomrule
\multicolumn{4}{@{}p{0.98\linewidth}@{}}{注：日付はすべて 2026 年（月-日）．コミットは git の短いハッシュで，登録は登録の文が入ったコミット（案と確定が別のものは確定の方），結果は結果を記録したコミット．「事前」「事前登録外」の行は，同じ日付のものも含めて，登録のコミットが結果のコミットより前である．「事前登録外」は，判定の形を登録せず，読み方や測り方だけを計算の前に記録したものである．} \\
\end{tabular}
